% RAY (Resilience Assistant for You) — OASIS @ ACM SIGSPATIAL 2026, Track A short paper.
% Build: pdflatex -> bibtex -> pdflatex x2 (see build.sh).
%
% OWNER DECISIONS PENDING (do not submit before resolving):
%   1. Author list and order — confirmed 2026-09-03 (Yang, Zou, Gong, Wang);
%      Ziyi Wang's address is from a public homepage, verify.
%   2. Page limit and template variant are assumed (sigconf, 4 pages +
%      references) — reconcile with the Track A brief behind the event
%      login before submission.
\documentclass[sigconf,screen]{acmart}

\setcopyright{rightsretained}
\copyrightyear{2026}
\acmYear{2026}
\acmConference[OASIS '26]{OASIS Challenge @ ACM SIGSPATIAL}{November 2026}{Minneapolis, MN, USA}
\acmBooktitle{OASIS Challenge at ACM SIGSPATIAL 2026, Track A: Disaster Resilience \& Vulnerability Analysis}
\settopmatter{printacmref=false}

\begin{document}

\title{RAY: Resilience Assistant for You}
\subtitle{An Accountable GeoAI System for Place-Based Disaster Intelligence}

\author{Yifan Yang}
\affiliation{%
  \institution{Department of Geography, Texas A\&M University}
  \city{College Station}\state{TX}\country{USA}}
\email{yyang295@tamu.edu}

\author{Lei Zou}
\affiliation{%
  \institution{Department of Geography, Texas A\&M University}
  \city{College Station}\state{TX}\country{USA}}
\email{lzou@tamu.edu}

\author{Wenjing Gong}
\affiliation{%
  \institution{Department of Landscape Architecture and Urban Planning, Texas A\&M University}
  \city{College Station}\state{TX}\country{USA}}
\email{wenjinggong@tamu.edu}

\author{Ziyi Wang}
\affiliation{%
  \institution{Department of Computer Science and Engineering, Texas A\&M University}
  \city{College Station}\state{TX}\country{USA}}
\email{ziyiwang@tamu.edu} % from the author's homepage (zoe-wang.com); verify before submission

\begin{abstract}
Most GeoAI disaster systems demonstrate autonomy: a language model that
runs a spatial pipeline. RAY (Resilience Assistant for You) demonstrates
\emph{accountability}: a risk-analyst agent that operates inside the
\emph{Steward Harness}, a technical layer that enforces outcome, process,
and institutional validity during operation rather than promising them in
prose. Every artifact is hashed into an append-only audit log; a
declarative claim policy scopes what the agent may assert by role,
evidence tier, spatial resolution, and geographic authority; a separate
distribution policy scopes what a build may publish, verified by a CI
gate that fails the deploy on violation; every factual sentence must cite
an artifact identifier or be refused; and a \emph{verifiability} axis
(retained $>$ re-derivable $>$ cited-only) states what a reader can do to
check a claim. We exercise the harness on a nationwide hourly hazard
watch and three deep cases across two hazard types --- the 2025 Eaton
Fire, Hurricane Milton (2024), and Hurricane Ian (2022) --- delivered
through an installable, keyless web app with resident and planner modes.
We then ask what it takes for a vision--language model's output to count
as evidence. Six open VLMs, served locally on one 24\,GB GPU, graded
1{,}225 street-view images and image pairs per model across the three
cases under verbatim published prompts; every prediction is a hashed
record, thresholds fail closed, and the resulting grids are flagged
\emph{model-derived}: publishable as an evaluation, unusable as a claim.
Single-image grading comes within two points of closed commercial
models; pre/post change grading does not, and each model fails with a
different fixed class bias. The test suite doubles as a verifiable
evaluation environment (363 Python and 67 app tests in CI), and the
audit record preserves three incidents where the harness caught the
system violating its own rules --- evidence, we argue, that the mechanism
does work that documentation alone does not.
\end{abstract}

\begin{CCSXML}
<ccs2012>
<concept>
<concept_id>10010147.10010178</concept_id>
<concept_desc>Computing methodologies~Artificial intelligence</concept_desc>
<concept_significance>500</concept_significance>
</concept>
<concept>
<concept_id>10002951.10003227.10003236</concept_id>
<concept_desc>Information systems~Geographic information systems</concept_desc>
<concept_significance>500</concept_significance>
</concept>
</ccs2012>
\end{CCSXML}
\ccsdesc[500]{Information systems~Geographic information systems}
\ccsdesc[500]{Computing methodologies~Artificial intelligence}

\keywords{GeoAI agents, accountability, disaster resilience, social
vulnerability, provenance, policy enforcement, vision--language models, H3}

\maketitle

\section{Introduction}

Large-language-model agents can now operate geographic information
systems: they select tools, write spatial queries, and produce maps and
narratives with little human intervention
\cite{li2023autonomous,zhang2024geogpt}. In disaster response this
capability meets an unforgiving audience. An emergency manager acting on
a tile-level damage estimate, or a resident reading a sentence about
their own street, cannot audit a transformer. The prevailing answer is
disclaimers --- operational platforms ship forecasts with the caveat that
they ``complement, do not replace'' official warnings --- which locates
accountability in the reader rather than in the system.

RAY locates it in the system. Following the framing of
\citet{yang2026accountable}, we treat an agent's claim as valid only
when three conditions hold at once: the \emph{outcome} passes executable
spatial checks, the \emph{process} that produced it is preserved and
inspectable, and the claim is \emph{institutionally} authorized --- the
right actor asserting the right thing at the right resolution. The
\emph{Steward Harness} enforces all three in code, and adds a fourth
question that arises the moment third-party data enters: what can a
\emph{reader} do to check this, and may the evidence be kept at all?

This paper describes the deployed system\footnote{Live site, app, and
repository: \url{https://rayford295.github.io/ray-resilience/}} and reports
what its enforcement actually caught. Our contributions are:
(1)~a harness architecture with two declarative policy planes --- claim
and distribution --- that fail closed, plus a per-claim
\emph{verifiability} axis with weakest-link semantics;
(2)~three harness-built deep cases across two hazard types, joining
structure-level damage ground truth, debris volumes, CDC social
vulnerability, and cross-view street imagery on H3 r9 grids, behind a
nationwide hourly watch;
(3)~a treatment of model output as \emph{governed evidence}: six open
vision--language models evaluated on 1{,}225 labelled samples per model
across the three cases, with every prediction recorded, thresholds that
fail closed, and results that the agent is forbidden to cite;
(4)~a test suite designed as a verifiable evaluation environment, and an
append-only audit record that preserves three self-caught violations ---
including one where the deployed site violated the project's own written
distribution constraint and the fix was a new policy plane, not a patch.

\section{The Steward Harness}

The harness sits between every pipeline stage, and between the language
model and every sentence it emits.

\textbf{Outcome validity.} Executable spatial checks run inside each
build: CRS assertions, raster$\times$vector and tract-join integrity,
sanity bounds, and \emph{mandatory} per-tile uncertainty fields --- a
grid cell without a declared uncertainty structure fails the build.

\textbf{Process validity.} Every registered artifact receives a SHA-256
digest and an append-only audit row naming its agent, inputs, and UTC
timestamp. Failures are recorded, never erased; the app renders lineage
from any map layer back to timestamped source snapshots. The gateway's
own audit is redacted by one rule --- it holds no finer location than the
claim plane lets any answer use: a point is stored as its H3 r9 cell, an
area's corners rounded to $\sim$110\,m, and the question as sha256 plus
length, so a caller can prove their question produced a row without the
log retaining what a resident typed about their own home.

\textbf{Institutional validity.} A single YAML policy file carries two
planes of ordered, first-match rules with default deny, validated at
construction time --- an unknown match key fails loudly at load rather
than silently broadening authorization. The \emph{claim plane} scopes
what may be asserted, by role (resident vs.\ planner), evidence tier,
spatial resolution, and geography: tile-level evidence never yields
parcel-level claims. The \emph{distribution plane} scopes what a build
may \emph{publish}, by artifact class, resolution cap, and audience; CI
verifies the assembled public site against it and fails the deploy on
violation. One \texttt{published\_events} list feeds both planes, so the
agent cannot read what the publisher withholds.
Section~\ref{sec:incidents} explains why the second plane and that
shared list exist.

\textbf{Citation by default.} Every factual sentence the agent produces
must cite an artifact identifier. The exemption set is closed
(questions, imperative safety advice, statements about what the answer
cannot say), so a sentence form nobody anticipated costs a refusal, not
an uncited claim. Refusals name the policy rule that triggered them.

\textbf{Verifiability.} Evidence tiers describe depth; they do not say
whether a reader can \emph{check} a claim. Each artifact therefore
carries \texttt{verifiability} $\in$ \{\texttt{retained},
\texttt{re-derivable}, \texttt{cited-only}\} with weakest-link
aggregation, and a \texttt{license} attribute gates redistribution by
rule. The axis earns its keep where retention is forbidden: map-platform
terms disallow storing map content, which is structurally incompatible
with proving traceability by hashing inputs. The harness resolves this
by attesting to the \emph{request} plus a response digest --- a lookup
record that is publishable precisely because it holds no content --- and
a containment property test enforces that no response payload ever
reaches the audit record. Grounded free-text generation is not
reproducible even at temperature zero, so it falls through to
default-deny: the citable unit is the citation, never the prose. A
fourth value, \texttt{open-license-attribution}, is exercised in
production for OSM-derived facility context (ODbL): redistributable, but
only with its attribution, which therefore travels inside the artifact.

\section{System and Deep Cases}

RAY is three loosely coupled planes: a data plane (GitHub Actions
pipelines through the harness, publishing hashed artifacts), a
presentation plane (an installable React + MapLibre PWA, keyless and
offline-cacheable), and an agent plane (any OpenAI-compatible endpoint
wrapped in policy pre-check $\rightarrow$ grounded generation
$\rightarrow$ claim post-check $\rightarrow$ audit). If the agent plane
is down, the maps and analysis keep working --- graceful degradation as
fail-closed design made visible.

\textbf{Tier 1 --- nationwide watch.} Four keyless federal connectors
(USGS earthquakes, NWS alerts, NHC tropical cyclones, NIFC/WFIGS
wildfires) refresh hourly in CI, writing append-only gzip snapshots and
a national watch product. Per-source failures are declared, not hidden:
the map badge reports \emph{mapped-of-total} features and names sources
that failed, and an ArcGIS error envelope fails the wildfire connector
closed instead of reporting zero fires. NWS alerts render hollow ---
an advisory about what may come never looks like an occurrence. A
Day-1 flash-flood outlook (NOAA/WPC, public domain) publishes as a
\emph{separate} product beside the watch, carrying its own declared
boundary: outlook only, no damage conclusions, never a replacement for
official warnings.

\textbf{Tier 2/3 --- three deep cases, one harness.} Exposure,
vulnerability, and damage analysis exist only inside three areas of
interest, and the app says so rather than extrapolating:
\begin{itemize}
\item \emph{Eaton Fire 2025 (CA):} 18{,}428 CAL FIRE DINS structure
  points $\rightarrow$ a 265-cell H3 r9 damage grid with per-severity
  counts and mandatory uncertainty; CDC SVI 2022 joined across 20 Census
  tracts with the downscaling approximation declared per feature;
  2{,}244 reliability-gated cross-view street-imagery samples
  $\rightarrow$ a 109-cell evidence grid; 46{,}341 residents (2020
  census blocks, centroid-allocated with the pre-event vintage declared)
  and 27 OSM critical facilities. A damage class with $n{=}30$
  is declared as lacking statistical power rather than hidden.
\item \emph{Hurricane Milton 2024 (FL):} 2{,}556 labeled pre/post
  street-view pairs $\rightarrow$ a 15-cell grid at Horseshoe Beach,
  each feature carrying the attribution caveat that post-imagery is
  season-cumulative (Debby, Helene, Milton); 5{,}618 Pinellas County
  cells with county debris-program volumes and wind/rain covariates.
  772{,}293 residents and 200 OSM facilities across the two AOIs. A
  generated-imagery dataset was excluded, and the exclusion itself is
  auditable in the event dossier.
\item \emph{Hurricane Ian 2022 (FL):} 886 matched samples $\rightarrow$
  a 190-cell evidence grid; 4{,}121 further street-view positions ship
  as a \emph{density-only} layer because no verifiable per-point
  severity link exists --- candor over coverage; 5{,}428 residents and
  79 OSM facilities.
\end{itemize}
All source imagery traces to a hashed dataset registry (134{,}272
files, $\sim$33\,GB, SHA-256). Facility context declares OSM
\emph{presence, never operational status}; population outside the
evaluated tiles is declared as a total, never mapped into tiles the
event has no evidence for; and no claim-plane rule authorizes facility
claims, so the agent's answers about them default-deny.

\textbf{Two modes.} Residents geocode an address and receive a
plain-language dossier --- stat callouts with their qualifiers attached,
declared unknowns rendered as prominently as findings, and one of three
coverage answers: covered, outside the evaluated areas, or
\emph{not determined} when a layer could not be read (the system will
not claim an address is outside areas it failed to look at). Planners
weight structural damage against social vulnerability on a live slider
(every adjustment audit-logged), inspect lineage, and can draw an area
and ask about it; the answer covers only the evaluated tiles the
selection contains, never merges statistics across events, and
highlights the tiles it cited.

\section{Model Output as Governed Evidence}
\label{sec:vlm}

The obvious next step for any imagery-rich disaster system is to let a
vision--language model (VLM) grade damage. The harness turns that step
into a question: under what conditions may a model's guess become
evidence? Our answer is mechanical. Model output is a distinct evidence
class, \texttt{model\_derived}, that carries its own measured accuracy
alongside it or carries nothing; no claim-plane rule authorizes a claim
from it; and the agent's evidence store refuses to load it. The
numbers below are therefore published as an \emph{evaluation of a
method on labelled cases}, never as facts about places.

\textbf{Setup.} We ported the Damage Recognition prompts of
RAPID~\cite{yang2026rapid} verbatim --- Prompt~C for a single post-event
image on the five CAL FIRE DINS classes, Prompt~B for a pre/post pair on
three human-perception classes --- and froze each as a
\texttt{prompt\_sha256} in every record, so a later edit is visible as a
different prompt rather than a silent change in the numbers. The model is
any OpenAI-compatible vision endpoint; here, open weights served by
Ollama on one RTX~3090 (24\,GB), so the run is keyless and repeatable
from a workstation. Three cases: \emph{Palisades 2025} (RAPID Dataset
C2: 295 street-view images, folder-level DINS truth, all graded);
\emph{Milton 2024} (the Bi-Temporal pair set, 100 pairs per class,
seed 2026); \emph{Eaton 2025} (the matched cross-view set, 300 per class
with the $n{=}30$ minority class fully included; Prompt~C sees the
post-event field image only --- declared, not hidden --- and its five
classes are collapsed onto the set's three repairability classes by DINS
semantics before scoring). Every prediction is a record: image sha256,
truth, prediction, response digest, latency, model digest. An
unparseable answer is recorded as unparseable, never coerced into a
class. A stage aborts if more than 20\,\% of answers are off-schema, if
a metric leaves $[0,1]$, if fewer than 95\,\% of graded samples fall
inside the grid the event already describes, or if any product lacks an
uncertainty block. We report accuracy and RAPID's normalized
cross-severity error, $\mathrm{NCSE}=\overline{|y-\hat{y}|}/(K{-}1)$.

\begin{table}[t]
\caption{Six open VLMs, one pass at temperature 0, quantised weights, on
the same seeded samples under the same prompt digests. Accuracy / NCSE.
Blocks are not comparable across columns (different class counts). The
closed-model row is RAPID's published number on the same Palisades
images and on a 150-pair subset of the Milton set.}
\label{tab:vlm}
\small
\setlength{\tabcolsep}{4pt}
\begin{tabular}{@{}lccc@{}}
\toprule
model & Palisades (5 cl.) & Milton pairs (3 cl.) & Eaton (3 cl.) \\
\midrule
qwen2.5vl:7b (ref.)      & 0.475 / 0.213 & \textbf{0.597} / 0.238 & 0.910 / 0.049 \\
qwen3-vl:32b$^{\dagger}$ & \textbf{0.554} / 0.176 & 0.381 / 0.431 & 0.901 / 0.053 \\
qwen2.5vl:32b$^{\dagger\ddagger}$ & 0.507 / 0.206 & 0.560 / 0.240 & 0.931 / 0.039 \\
gemma3:27b               & 0.529 / 0.187 & 0.500 / 0.257 & 0.914 / 0.049 \\
mistral-small3.2:24b     & 0.505 / 0.203 & 0.487 / 0.258 & \textbf{0.935} / 0.037 \\
gemma3:12b               & 0.546 / 0.178 & 0.427 / 0.287 & 0.906 / 0.051 \\
qwen3-vl:8b              & 0.514 / 0.195 & 0.462 / 0.324 & 0.924 / 0.041 \\
\midrule
closed (RAPID)           & 0.573 GPT-5-mini & 0.591 GPT-5.1 & --- \\
\bottomrule
\end{tabular}
\par\smallskip\footnotesize
$^{\dagger}$8k context so the weights fit the card; $^{\ddagger}$95\,\% of
the model in VRAM, the rest on CPU --- recorded in the run, because a
spilled model is a different run. A seventh model (llama3.2-vision:11b)
could not be loaded by the serving stack and is declared as a serving
failure, not scored.
\end{table}

\textbf{What the runs say.} Table~\ref{tab:vlm} reports seven models
(a 7B reference plus six candidates), each committed with its own eval
summary, prediction records, agreement grid, and GPU placement as it
finished. Four things hold across the table.
(i)~\emph{On a single post-event image, every open model beats the 7B
reference, and the best is within two points of the paper's closed
models.} Size is not the axis: gemma3:12b (0.546) is within one point
of qwen3-vl:32b (0.554) and ahead of gemma3:27b.
(ii)~\emph{On pre/post pairs, no open model beats the 7B reference, and
the failures are systematic.} qwen3-vl:32b calls 258 of 300 pairs
Severe (Mild recall 0.01); gemma3:12b calls 262 Moderate (Severe recall
0.17); mistral puts 213 in Moderate (Mild recall 0.13). Each model has a
different fixed bias on the change task, so a vote across them would not
be a fix; the per-class recall column in the committed table is the
finding.
(iii)~\emph{Eaton's small class is where every model is weak.}
\texttt{damaged\_repairable} ($n{=}30$) recall runs from 0.03 to 0.43
while the two large classes sit above 0.9; the headline 0.93 is the
large classes.
(iv)~\emph{Latency is a model property.} The qwen3 family spends its
time in hidden reasoning before a 60-character answer: qwen3-vl:8b's
median is 5\,s, but nine Palisades images took over 30\,s, and on Milton
24 pairs hit the client timeout and were recorded as unavailable
(8.3\,\% unanswered, under the 20\,\% bound, so the case stands with its
unanswered rate on the record).

\textbf{What the harness did with it.} The agreement grids sit beside
the human-labelled grids they are evaluated against, flagged
\texttt{model\_derived}; the publication planner classifies them as
evaluation products; the claim plane has no rule that admits them; and
the gateway's evidence store skips them and records the exclusion. A
reader can verify every number from the committed prediction records
and prompt digests without a GPU. What no reader can do --- and what the
system therefore does not do --- is turn a 0.93 into a sentence about a
street.

\section{The Harness Catching Its Own System}
\label{sec:incidents}

Three preserved incidents are, we believe, the strongest evidence the
mechanism does real work.

\textbf{The publication boundary.} A parcel-level DINS file (18{,}428
points with per-structure damage) reached the public site through an
asset-copying script, violating four written statements in the
repository --- while breaking zero policy rules, because the claim plane
governs what the agent \emph{asserts} and nothing had asserted the file.
The fix was structural: a distribution plane in the same policy file,
a publication-boundary planner, and a CI gate that verifies the
assembled site; the public surface dropped from 30 files to 16. The
incident report is committed alongside the code it indicts.

\textbf{Citation by default, inverted.} The claim post-check originally
required citations only on sentences containing digits, so
``\emph{Your neighborhood was not significantly affected}'' passed
uncited. The check was inverted to citation-by-default with a closed
exemption set; acceptance on a representative draft set fell from 9/13
to 7/13 sentences --- the two newly refused sentences being exactly the
uncited reassurances a resident should not receive.

\textbf{The agent reading what the publisher withheld.} The gateway's
evidence store loaded every event dossier and grid on disk, ignoring
both the \texttt{published\_events} list and the \texttt{model\_derived}
flag. The first evaluation case written under \texttt{events/} --- the
Palisades VLM run of Section~\ref{sec:vlm} --- would have been located by
point and its predictions cited as tile facts, and Milton's model-graded
pair grid would have been served beside the human-labelled grid it is
evaluated against. The publication planner had excluded them; the agent
had not. The fix makes both planes publish from one list, has the store
skip model-derived products and record every exclusion, and refuses to
start a store over ``every event on disk'' as a misconfiguration rather
than a fallback (20 new tests). The pattern is the same as the first
incident: a constraint written in one place and enforced in another is
enforced nowhere.

\section{Evaluation Environment}

Following \citet{yang2026accountable}, the test suite is designed as the
paper's verifiable evaluation environment rather than an afterthought:
363 Python tests and 67 app tests run in CI on every push, including a
cell-by-cell policy-matrix sweep (every role $\times$ tier $\times$
resolution combination against expected authorization), adversarial
claim-check cases, fail-closed connector tests against real error
envelopes, the containment property test for content-free lookup
records, the gateway-scope tests above, and the VLM stages' parsing,
fail-closed and resume logic with the model call injected. Anchor checks
resolve every path-shaped reference in the documentation against the
repository, so the manual cannot silently rot. A CI check regenerates
the publication allowlist from the policy and fails on any drift, which
is how a hand edit cannot widen the public surface. Judges can run the
environment in two commands from the README.

\section{Related Work}

Autonomous-GIS agents demonstrate that LLMs can operate spatial tooling
\cite{li2023autonomous,zhang2024geogpt}; our contribution is orthogonal
--- a harness that constrains any such agent's claims to its evidence.
RAPID~\cite{yang2026rapid} shows that a multi-agent pipeline over closed
APIs can grade damage from imagery interpretably; we keep its prompts,
metric and acceptance rules as harness stages, replace the closed models
with open ones, and decline its LLM task planner, because a model
choosing which agents run conflicts with a policy plane. Operational
platforms such as the Texas Disaster Information System~\cite{tdis2026}
field forecast-forward impact summaries at state scale; accountability
there rests on interface disclaimers, whereas RAY binds each sentence to
a hashed artifact and a policy rule. Global registries
(EM-DAT~\cite{delforge2025emdat}, GDACS~\cite{gdacs}) record that events
happened; none carries, per row, what resolution of claim the evidence
authorizes or whether a reader may keep it --- the two fields this system
computes for every artifact.

\section{Limitations}

The deep-case builders read a $\sim$33\,GB corpus on the maintainer's
workstation, so third parties can verify the committed artifacts,
hashes, and audit logs but not regenerate them. The VLM results are one
pass at temperature zero on 4-bit quantised weights; a model digest pins
the weights, not the arithmetic, and 32B is the ceiling of a 24\,GB card.
Milton and Eaton are seeded stratified samples, and the closed-model
reference for Milton was graded on a different 150-pair subset. The agent
gateway runs locally (no hosted endpoint; no auth or rate limiting yet),
so the public demo ships without chat. The re-derivable regime is
implemented and tested against a stub but has never run against a live
commercial API. Nationwide coverage is Tier-1 watch only; the app
declares this boundary instead of smoothing over it, which we consider a
feature, but reviewers should not mistake three deep cases for national
analysis capability.

\begin{acks}
RAY is an MIT-licensed research prototype. It is not an official
forecasting or warning service.
\end{acks}

\bibliographystyle{ACM-Reference-Format}
\bibliography{references}

\end{document}
